\section{Conservative reset models and large irreducibles}\label{app:reset-models}
\subsection{Two equations and the conservative models}

An algebra $(A,\del)$ is \emph{conservative} if $x\del y\in\{x,y\}$.
The three-element generator is conservative, although winner-set reset
and free reset algebras generally are not. Here we classify conservative
models of the equational theory; their elements need not themselves be
concrete winner sets.
The order-algebra account of this variety already yields the forest
restriction \cite[Theorem~3.3]{FreeseEtAl2002}. Our convention reverses
that order. The proof below records the short two-law reconstruction
inside the conservative class and supports its mechanization.

A \emph{forest order} is a partial order in which the upper set of each
element is linearly ordered. In the finite case its Hasse diagram is a
rooted forest, with roots above their descendants. Define
\begin{equation}\label{reset:eq:forest-op}
 x\del_P y=\begin{cases}x,&y\leq_P x,\\y,&\text{otherwise}.\end{cases}
\end{equation}
Comparable elements yield the higher one; incomparable elements yield
the right argument.

\begin{theorem}[Conservative presentation]\label{reset:thm:forest}
For a conservative algebra, the following are equivalent:
\begin{enumerate}
\item it satisfies every reset identity;
\item it satisfies the two equations
\begin{align}
 (x\del y)\del x&=y\del x,\label{reset:eq:absorb}\\
 x\del(y\del z)&=y\del(x\del z);\label{reset:eq:left-comm}
\end{align}
\item it has form \eqref{reset:eq:forest-op} for a unique forest order.
\end{enumerate}
The order is recovered by
$x\leq_P y$ iff $x\del y=y\del x=y$.
\end{theorem}
\begin{proof}
Both equations hold for winner-set reset. In
\eqref{reset:eq:absorb}, a resolved $x$ returns $x$ on both sides; an
ambiguous $x$ gives $x\cup y$ on both sides.
In \eqref{reset:eq:left-comm}, a resolved $z$ returns $z$ on both sides;
an ambiguous $z$ gives $x\cup y\cup z$.

Assume the two equations and conservativity, and use the proposed
order. Reflexivity follows from conservativity, and antisymmetry is
immediate. If $x\del y=x$, absorption implies $y\del x=x$.
For transitivity suppose $x\leq y\leq z$. Left commutation gives
\[
 z=z\del(y\del x)=y\del(z\del x).
\]
If $z\del x=z$, the preceding absorption argument gives $x\del z=z$.
Otherwise conservativity gives $z\del x=x$, and the display gives
$z=y$, reducing to the assumed relation $x\leq y$.
If $x\leq y$ and $x\leq z$, left commutation gives
$y\del z=z\del y$; conservativity then makes $y,z$ comparable.
The operation is recovered from the order because selecting the left
argument forces the two-sided comparison just proved.

Conversely, \eqref{reset:eq:forest-op} satisfies both equations.
For left commutation, if $z$ is below both $x,y$, these ancestors
are comparable and both sides return the higher one. If it is below
only one, both sides return that ancestor; otherwise they return $z$.
Absorption follows directly from the comparable and incomparable cases.
This order reconstruction in both directions is kernel-checked in
\texttt{ResetModels.lean}. Lemma~\ref{reset:lem:forest-membership} below proves
that every forest operation satisfies all reset identities, completing
the equivalence.
\end{proof}

The qualification ``conservative'' is essential: the two equations are
a complete presentation within that class, not a proposed equational
basis over arbitrary algebras.

\begin{lemma}[Closure constructions]\label{reset:lem:closures}
The reset variety is closed under the following constructions:
\begin{enumerate}
\item adjoin one new element absorbing on both sides;
\item take a disjoint union of models, using the component operation
within a component and returning the right argument across components.
\end{enumerate}
\end{lemma}
\begin{proof}
Every reset identity is regular: the two-element ambiguous subalgebra
$\{\{a,b\},\{a,b,c\}\}$ is a semilattice under reset, and its identities
have equal variable supports. This is the standard semilattice argument
for regular identities, in the setting of P{\l}onka sums
\cite{Plonka1967}.
Thus adjoining an absorbing element preserves every identity: either
both sides encounter it or neither does.

For the disjoint union, an evaluated term belongs to the component of
its rightmost variable. Both sides of a reset identity have the same
rightmost variable, by their behavior on atomic valuations. Fix the
set $B$ of input positions whose values belong to that component.
Prune a term with root in $B$ as follows. Retain its right child and
prune recursively; retain its left child exactly when that child's
root is in $B$. The result uses variables in $B$ and evaluates the
original term in the chosen component.

For every $C\subseteq B$, its selector is $E_t(C)\cap B$.
To see this by induction, note that $\rho(s)\notin C$ implies
$E_s(C)=\{\rho(s)\}$. Thus a removed left child contributes no
position in $B$ to the original selector, and a retained one obeys the
selector recursion. Equal original selectors give equal pruned selectors,
so the pruned terms form a reset identity. The component model satisfies
that identity. This proves closure.
\end{proof}

\begin{lemma}[Forest membership]\label{reset:lem:forest-membership}
Every forest-order operation belongs to the reset variety.
\end{lemma}
\begin{proof}
Every finite rooted forest is constructed from singletons by disjoint
union and by adjoining an absorbing root to each component. These are
exactly the operations of Lemma~\ref{reset:lem:closures}.
For an infinite forest order, any finite set of elements is a subalgebra,
by conservativity, and its induced order is a finite forest order.
Every evaluation of an identity lies in such a finite subalgebra.
\end{proof}

As a counting consequence, there are $(n+1)^{n-1}$ conservative reset
algebras on a fixed labeled $n$-element carrier. Adjoining a new common
root gives a bijection from rooted forests on that carrier to labeled
trees on $n+1$ vertices, so the formula is the classical tree count.
The first five counts are $1,3,16,125,1296$.

\subsection{Large irreducible models from combs}

Three-element completeness concerns identities. It does not bound the
irreducible models of those identities. The forest description makes
this distinction explicit.

Let $L$ be any linear order with least element $0$. Its comb has
elements
\[
 a,\qquad b_i, p_i\quad(i\in L).
\]
The $p_i$ form the ascending spine; $b_i$ is a leaf below $p_i$,
and $a$ is below every spine element. Precisely,
\[
 p_i\leq p_j\Longleftrightarrow i\leq j,\qquad
 b_i\leq p_j\Longleftrightarrow i\leq j,\qquad a\leq p_j.
\]
All other distinct pairs are incomparable. This is a forest order.
The leaves act by right projection, and its only other products are
\[
 p_i\del a=p_i,\quad
 p_i\del b_j=
 \begin{cases}p_i,&j\leq i,\\b_j,&j>i,\end{cases}\qquad
 p_i\del p_j=p_{\max(i,j)}.
\]

\begin{theorem}[Comb monolith]\label{reset:thm:comb}
The comb is subdirectly irreducible. Its least nonidentity congruence
has the single nonsingleton class $\{a,b_0\}$.
\end{theorem}
\begin{proof}
The map replacing $a$ by $b_0$ and fixing every other element is an
endomorphism: every spine has the same product with $a$ and $b_0$.
Its kernel is the stated two-point congruence.

Any congruence identifying a spine $p_i$ with a leaf identifies $a,b_0$:
multiply the pair on the right first by $a$, then by $b_0$.
The spine returns $p_i$ in both cases, while the leaf returns $a,b_0$.
It remains to reduce each other distinct pair to a spine--leaf pair.
\begin{itemize}
\item For $a\sim b_j$ with $j>0$, multiply on the left by $p_0$ to
obtain $p_0\sim b_j$. For $j=0$ the desired pair is already identified.
\item For $b_i\sim b_j$ with $i<j$, multiply on the left by $p_i$ to
obtain $p_i\sim b_j$.
\item For $p_i\sim p_j$ with $i<j$, multiply on the right by $b_j$
to obtain $b_j\sim p_j$.
\end{itemize}
Every nonidentity congruence therefore contains the displayed kernel.
The argument, including the endomorphism and leastness, is
kernel-checked for arbitrary $L$ in \texttt{ResetModels.lean}.
\end{proof}

\begin{corollary}[Known residual largeness, explicit comb witnesses]\label{reset:cor:residual}
The reset variety, equivalently $\operatorname{V}(G_3^{\rm op})$, has
subdirectly irreducible models of size $2n+1$ for every $n\geq1$, and
of every infinite cardinality. In particular, it has no cardinal bound
on the sizes of its subdirectly irreducible members.
\end{corollary}
\begin{proof}
Use an $n$-element chain for the finite examples. For an infinite
cardinal $\kappa$, use its initial ordinal as $L$; the comb has size
$1+2\kappa=\kappa$. Lemma~\ref{reset:lem:forest-membership} places every comb
in the variety, and Theorem~\ref{reset:thm:comb} proves irreducibility.
\end{proof}

Residual largeness of this variety was already stated by Freese et al.
\cite[Section~3, before Theorem~3.2]{FreeseEtAl2002}. The construction
above gives explicit comb witnesses and a generic formal monolith proof;
priority for this particular witness construction is unestablished.

Residual largeness is compatible with the finite basis in
Theorem~\ref{reset:thm:basis}. Moreover, the semilattice-based Willard route
was unavailable for reset independently of this result: two distinct
atomic singletons form a right-projection subalgebra, whose term
operations are projections, so the reset variety has no Taylor term.
These facts concern the reset reduct; adding the other owned-choice
operations changes subalgebras and congruences. The full owned-choice
algebra has semilattice operations, so it does not inherit this
obstruction or the residual conclusion automatically.
